% Three scales kept apart, and this chapter's evidence marking.
\begin{tikzpicture}[font=\sffamily\scriptsize, x=1cm, y=1cm]

\node[chlabel, anchor=west] at (-6.8,3.4) {three scales that arrive tangled in almost every requirement document};

\node[regcell, fill=white, minimum width=126mm, anchor=north west,
      align=left, inner sep=3pt] at (-6.8,3.0)
  {\begin{tabular}{@{}p{26mm}p{30mm}p{34mm}p{22mm}@{}}
   \textbf{axis} & \textbf{document} & \textbf{what it specifies} &
   \textbf{values} \\[1pt]
   \hline
   stop \textbf{category} & the electrical equipment standard, clause 9.2.2 &
   what the stopping action does to actuator power & \textbf{0, 1, 2} \\[2pt]
   \textbf{performance} level & the machinery control parts standard &
   how \textbf{reliably} the function is performed & \textbf{a to e} \\[2pt]
   \textbf{integrity} level & the machinery functional safety standard, and the
   general one & the same question, in a different sector's vocabulary &
   \textbf{1 to 3}, and 4 in the general one \\
   \end{tabular}};

\node[regcell, fill=usrcol, minimum width=126mm, text width=124mm, inner sep=3pt,
      anchor=north west, align=left] at (-6.8,0.2)
  {\textbf{A well formed requirement names one value from the first row and one
   from the second or third:} "emergency stop shall be stop category 1, achieved
   with performance level d". \textbf{There is no stop category 3 and no
   performance level 1.} A category 0 stop implemented with a single undiagnosed
   contactor is still a category 0 stop, at a poor performance level, because the
   two axes are independent of each other.};

\node[regcell, fill=red!18, minimum width=126mm, text width=124mm, inner sep=3pt,
      anchor=north west, align=left] at (-6.8,-1.7)
  {\textbf{This volume prints no cross-map between the two reliability scales.}
   Such tables circulate widely, they are a convenience rather than a normative
   equivalence, and no primary verified mapping was obtained during this volume's
   research. Printing one would be repeating something rather than knowing it.};

% ================================================================ evidence
\node[chlabel, anchor=west] at (-6.8,-3.4) {and this chapter's evidence marking, because not all of it is the same quality};

\node[regcell, fill=white, minimum width=126mm, anchor=north west,
      align=left, inner sep=3pt] at (-6.8,-3.8)
  {\begin{tabular}{@{}p{30mm}p{40mm}p{44mm}@{}}
   \textbf{class} & \textbf{what it means} & \textbf{what is in it here} \\[1pt]
   \hline
   \textbf{read directly} & opened and read in its own document, usually in a
   free preview & the safe state definition, the tool classification, the six
   emergency stop requirements, the complementary measure clause \\[2pt]
   \textbf{attested} & several independent sources agree and the wording is
   identical, but no primary document was opened & the three stop category
   definitions, and one clause that lies past the end of a preview \\[2pt]
   \textbf{practice} & \textbf{no citable source found at all} & that a safe
   state must be reachable from every other state, and that a fault latches \\
   \end{tabular}};

\node[umlnote, text width=126mm, anchor=north west] at (-6.8,-8.2)
  {The third row is the useful one. Two of the most important rules in this
   chapter have no source behind them, which is not a reason to leave them out
   and is every reason to say so. A volume that cites a document for a practice
   nobody has read it in is doing something worse than admitting the gap.};
\end{tikzpicture}
